\begin{frame}{Vì sao post rời rạc không đủ cho reward model}
  \centering
  \begin{tikzpicture}[
    node distance=0.75cm,
    every node/.style={font=\small},
    box/.style={draw, rounded corners, text width=2.5cm, minimum height=1cm, align=center}
  ]
    \node[box, fill=softblue] (pa) {Prompt A};
    \node[box, fill=softgreen, right=1cm of pa] (oa) {Post A\\score cao};
    \node[box, fill=softblue, below=1.2cm of pa] (pb) {Prompt B};
    \node[box, fill=softorange, right=1cm of pb] (ob) {Post B\\score thấp};
    \node[box, fill=softred, right=1.3cm of oa, yshift=-0.6cm, text width=3.2cm, minimum height=2.2cm] (wrong) {\textbf{Ghép sai}\\Post A = chosen\\Post B = rejected};

    \draw[-{Stealth}, thick] (pa) -- (oa);
    \draw[-{Stealth}, thick] (pb) -- (ob);
    \draw[-{Stealth}, thick, dashed] (oa) -- (wrong);
    \draw[-{Stealth}, thick, dashed] (ob) -- (wrong);
  \end{tikzpicture}

  \vspace{1em}
  \begin{alertblock}{Vấn đề lý thuyết}
    Chosen/rejected chỉ có ý nghĩa khi hai output cùng xuất phát từ một prompt. Ghép hai post khác prompt khiến reward model học tương quan giả.
  \end{alertblock}
\end{frame}
